% The four parts of the Day 14 lab. The steps follow Labs/day14/README.md.
% The budget values are the examples of the README: the values of Day 14,
% Part 3 of the lecture. budget.py printed them.

% ------------------------------------------------------------------ Part A
\begin{frame}{Part A: a team and a problem (20 min)}
  \small
  \begin{enumerate}
    \item \textbf{Form the team (5 min).} Join with a second lab group.
          Write the names and the two group names in \code{proposal.md}
          and \code{report.md}.
    \item \textbf{Select a problem (10 min).} Read Sections 2 and 7 of
          \code{Docs/capstone.md}. Take one example, change one, or write
          your own problem.
    \item \textbf{Name the riskiest assumption (5 min).} Which part has the
          largest chance to fail? Part D tests it.
  \end{enumerate}
  \vskip 0.3em
  \footnotesize
  \renewcommand{\arraystretch}{1.15}
  \setlength{\tabcolsep}{4pt}
  \begingroup\centering
  \begin{tabular}{@{}lL{0.30\textwidth}L{0.24\textwidth}L{0.26\textwidth}@{}}
    \toprule
    & \textbf{Example} & \textbf{XIAO} & \textbf{Raspberry Pi} \\
    \midrule
    A & Machine monitor & Vibration model, alarm & Persons, dashboard \\
    B & Voice-controlled camera & Keyword model & Detector for 30 s \\
    C & Room counter with privacy & Person or no person & Count of persons \\
    D & Offline lab assistant & Gesture model & Language model, notes \\
    \bottomrule
  \end{tabular}\par\endgroup
\end{frame}

\begin{frame}{Part A: the problem, not a solution}
  \footnotesize
  \renewcommand{\arraystretch}{1.2}
  \setlength{\tabcolsep}{4pt}
  \begingroup\centering
  \begin{tabular}{@{}L{0.44\textwidth}L{0.48\textwidth}@{}}
    \toprule
    \textbf{A solution (weak start)} & \textbf{A problem (good start)} \\
    \midrule
    ``Run YOLO on a camera'' & Nobody should stand near a machine that
      has an alarm \\
    ``Use a language model'' & Students need lab answers when the internet
      is down \\
    ``Classify sounds'' & A camera should start only when a person asks
      for it \\
    \bottomrule
  \end{tabular}\par\endgroup
  \vskip 0.6em
  \small
  \begin{itemize}
    \item The problem names a user, a place, and a benefit. Then the
          requirements follow from it.
    \item Keep the scope small: reuse the labs of Days 1 to 13 and add one
          new part.
    \item A riskiest assumption is concrete: ``the IMU separates a normal
          fan from a fan with a weight on one blade''.
  \end{itemize}
\end{frame}

% ------------------------------------------------------------------ Part B
\begin{frame}{Part B: write the proposal (50 min)}
  \begin{columns}[T]
    \begin{column}{0.54\textwidth}
      \small
      \textbf{Steps}
      \begin{enumerate}
        \item Problem and requirements (15 min)
        \item Boards, model, and data (15 min)
        \item Budgets with \code{budget.py} (15 min)
        \item The five requirements (5 min)
      \end{enumerate}
      \vskip 0.4em
      One page. Write numbers where the template asks for numbers. Mark
      each value that you did not measure as ``estimate''.
    \end{column}
    \begin{column}{0.42\textwidth}
      \begin{block}{Sections of \texttt{proposal.md}}
        \small
        \begin{itemize}
          \item Problem and user
          \item Requirements with numbers
          \item Boards and split of the work
          \item Model and optimization
          \item Data
          \item Local decision and telemetry
          \item Budgets
          \item The five requirements
          \item Riskiest assumption and first test
        \end{itemize}
      \end{block}
    \end{column}
  \end{columns}
\end{frame}

\begin{frame}{Part B: requirements with numbers}
  \footnotesize
  \renewcommand{\arraystretch}{1.15}
  \setlength{\tabcolsep}{4pt}
  \begingroup\centering
  \begin{tabular}{@{}lL{0.24\textwidth}L{0.30\textwidth}L{0.30\textwidth}@{}}
    \toprule
    \textbf{No.} & \textbf{Requirement} & \textbf{Number} & \textbf{Test}
      \\
    \midrule
    R1 & Classify the vibration & Recall of ``abnormal'' at least 0.9 & A
      session that the training does not see \\
    R2 & Alarm at the machine & Within 3 s of the start & 10 timed events
      \\
    R3 & Alarm with no network & Works with the router off & Switch off the
      Wi-Fi in the demonstration \\
    \bottomrule
  \end{tabular}\par\endgroup
  \vskip 0.6em
  \small
  \begin{itemize}
    \item Rows of the example proposal (machine monitor,
          \code{solutions/proposal\_example.md}).
    \item Each requirement has a number and a test that you can run before
          the demonstration on Day 15.
    \item ``Fast'', ``accurate'', and ``small'' are not requirements. Nobody
          can check them.
  \end{itemize}
\end{frame}

\begin{frame}[fragile]{Part B: the budgets with \texttt{budget.py}}
\begin{lstlisting}[style=shell, basicstyle=\ttfamily\scriptsize]
python3 budget.py fit --flash-kb 3264 --ram-kb 298.68 --program-kb 300 \
    --params 1534 --bits 8 --peak-values 1028
python3 budget.py latency window=2000 model=100 decision=400 mqtt=100 \
    --limit-ms 3000
python3 budget.py power --work-mw 217 --rest-mw 19 --work-ms 50 \
    --period-ms 2000 --radio-mw 290.4 --radio-s-per-hour 3 --days 30
\end{lstlisting}
  \footnotesize
  \renewcommand{\arraystretch}{1.1}
  \setlength{\tabcolsep}{4pt}
  \begin{tabular}{@{}lL{0.80\textwidth}@{}}
    \toprule
    \textbf{Command} & \textbf{You see (the values of the lecture)} \\
    \midrule
    \code{fit} & Flash 308\,734 of 3\,342\,336 bytes, RAM 1028 of
      305\,848 bytes: fits \\
    \code{latency} & Total 2600 ms, margin 400 ms (met) \\
    \code{power} & 24.192 mW, 6.4 days; 30 days need 17.42 Wh (4708 mAh)
      \\
    \bottomrule
  \end{tabular}
  \vskip 0.3em
  \small
  Two more commands: \code{data} (volume of a message stream) and
  \code{cost} (one site and all sites). Use your own values, and name the
  source of each input.
  \source{Example values of Day 14, Part 3. The program size of 300 KB is
    an estimate.}
\end{frame}

\begin{frame}{Part B: check the five requirements}
  \footnotesize
  \renewcommand{\arraystretch}{1.15}
  \setlength{\tabcolsep}{4pt}
  \begingroup\centering
  \begin{tabular}{@{}lL{0.38\textwidth}L{0.48\textwidth}@{}}
    \toprule
    \textbf{No.} & \textbf{Requirement} & \textbf{In the machine monitor}
      \\
    \midrule
    1 & XIAOML Kit and Raspberry Pi & The XIAO detects the vibration; the
      Pi sees the persons and runs the dashboard \\
    2 & One optimized model, size and accuracy before and after & The XIAO
      model in \code{float32} and \code{int8}, recall on the test session
      \\
    3 & One local decision, no internet & The alarm on the XIAO; test with
      the router off \\
    4 & Telemetry with MQTT to a dashboard & The broker on the Pi, the
      dashboard of Day 13 \\
    5 & Benchmark table: latency, memory, energy & Both models, with the
      method of Day 9 \\
    \bottomrule
  \end{tabular}\par\endgroup
  \vskip 0.6em
  \small
  If one requirement is not met, change the design now, not on Day 15.
\end{frame}

% ------------------------------------------------------------------ Part C
\begin{frame}{Part C: the design review (40 min)}
  \begin{columns}[T]
    \begin{column}{0.46\textwidth}
      \small
      \textbf{About 8 minutes for each team}
      \begin{enumerate}
        \item \textbf{Present (3 min):} the problem, the split of the work,
              the budgets, the riskiest assumption.
        \item \textbf{Answer (5 min):} the questions of the instructor,
              with numbers.
      \end{enumerate}
      \vskip 0.3em
      While you wait: prepare the answers, then prepare the first test of
      Part D.
    \end{column}
    \begin{column}{0.50\textwidth}
      \begin{block}{The seven questions}
        \footnotesize
        \begin{enumerate}
          \item A number and a test for each requirement?
          \item A value and a source for each budget?
          \item Does the model fit: flash, RAM, time?
          \item The data: source, size, labels, test set?
          \item No network, a failed sensor, a bad update?
          \item Which data leaves? Who can send a command?
          \item The riskiest assumption and its first test?
        \end{enumerate}
      \end{block}
    \end{column}
  \end{columns}
\end{frame}

\begin{frame}{Part C: the result of the review}
  \small
  \begin{itemize}
    \item The instructor fills \code{review\_sheet.md}: the five
          requirements, the seven questions, the scope, and the result.
    \item \textbf{You record} in \code{report.md}: the result, the changes,
          and the weak points.
    \item Make the changes in \code{proposal.md} before Part D.
  \end{itemize}
  \vskip 0.4em
  \begin{exampleblock}{Example result (machine monitor)}
    \small
    Approved with changes. (1) The test session of R1 is on a different day
    or position. (2) The benchmark table needs the measured latency of the
    XIAO model, not the budget value of 100 ms. Weak point: the battery
    gives 6.4 days, not 30; the team uses a USB cable.
  \end{exampleblock}
  \vskip 0.3em
  \keypoint{A weak point is a result of the review, not a failure of the
    team.}
\end{frame}

% ------------------------------------------------------------------ Part D
\begin{frame}{Part D: the first test (40 min)}
  \small
  \begin{enumerate}
    \item \textbf{Plan (5 min):} which board, which lab folder, which data,
          which number decides.
    \item \textbf{Run (30 min):} reuse the lab folder of the day that the
          test needs.
    \item \textbf{Decide and plan (5 min):} the result, the decision, the
          plan for Day 15.
  \end{enumerate}
  \vskip 0.3em
  \footnotesize
  \renewcommand{\arraystretch}{1.15}
  \setlength{\tabcolsep}{4pt}
  \begingroup\centering
  \begin{tabular}{@{}lL{0.38\textwidth}L{0.46\textwidth}@{}}
    \toprule
    & \textbf{Riskiest assumption} & \textbf{First test} \\
    \midrule
    A & The IMU separates the two states & Two sessions with the Day 2
      logger, the Day 3 notebook \\
    B & Few false starts in a noisy room & 5 minutes of normal speech with
      the Day 5 model \\
    C & The XIAO model works in the room light & 20 photos at two times of
      the day \\
    D & Short and fast answers & 10 questions, time and length (Day 10) \\
    \bottomrule
  \end{tabular}\par\endgroup
\end{frame}

\begin{frame}{Part D: decide, and plan Day 15}
  \small
  \textbf{You record:} the test, the result with numbers, the decision
  (continue or change), and the plan.
  \vskip 0.4em
  \footnotesize
  \renewcommand{\arraystretch}{1.15}
  \setlength{\tabcolsep}{4pt}
  \begin{tabular}{@{}lL{0.70\textwidth}r@{}}
    \toprule
    & \textbf{Example plan for Day 15 (machine monitor)} & \textbf{Time} \\
    \midrule
    1 & Record the test session, train the final model & 30 min \\
    2 & Convert to \code{int8}, deploy & 30 min \\
    3 & Connect the programs of Days 12 and 13 & 40 min \\
    4 & Measure the benchmark table & 40 min \\
    5 & Rehearse the demonstration with the router off & 20 min \\
    \bottomrule
  \end{tabular}
  \vskip 0.5em
  \small
  \begin{itemize}
    \item Day 15 has 100 minutes in the lecture time and 80 minutes in lab
          Part A for the build. Plan less than 180 minutes.
    \item Plan B for the demonstration: a recorded input or a video.
  \end{itemize}
\end{frame}
